\markdownRendererDocumentBegin
\markdownRendererSectionBegin
\markdownRendererHeadingOne{Teoría de Números}\markdownRendererInterblockSeparator
{}\markdownRendererSectionBegin
\markdownRendererHeadingTwo{Números Primos y Factores}\markdownRendererInterblockSeparator
{}\markdownRendererStrongEmphasis{Definición:}\markdownRendererHardLineBreak
{}Un número \markdownRendererDollarSign{}a\markdownRendererDollarSign{} es \markdownRendererStrongEmphasis{factor} (o \markdownRendererStrongEmphasis{divisor}) de \markdownRendererDollarSign{}b\markdownRendererDollarSign{} si \markdownRendererDollarSign{}a\markdownRendererDollarSign{} divide a \markdownRendererDollarSign{}b\markdownRendererDollarSign{}, denotado \markdownRendererDollarSign{}a \markdownRendererBackslash{}mid b\markdownRendererDollarSign{}.\markdownRendererHardLineBreak
{}Un entero \markdownRendererDollarSign{}n > 1\markdownRendererDollarSign{} es \markdownRendererStrongEmphasis{primo} si sus únicos divisores positivos son \markdownRendererDollarSign{}1\markdownRendererDollarSign{} y \markdownRendererDollarSign{}n\markdownRendererDollarSign{}.\markdownRendererParagraphSeparator
{}\markdownRendererStrongEmphasis{Factorización prima única:}\markdownRendererHardLineBreak
{}Todo entero \markdownRendererDollarSign{}n>1\markdownRendererDollarSign{} se puede escribir de forma única (hasta orden) como:\markdownRendererSoftLineBreak
{}\markdownRendererDollarSign{}\markdownRendererDollarSign{}\markdownRendererSoftLineBreak
{}n = p\markdownRendererUnderscore{}1\markdownRendererCircumflex{}\markdownRendererLeftBrace{}\markdownRendererBackslash{}alpha\markdownRendererUnderscore{}1\markdownRendererRightBrace{} p\markdownRendererUnderscore{}2\markdownRendererCircumflex{}\markdownRendererLeftBrace{}\markdownRendererBackslash{}alpha\markdownRendererUnderscore{}2\markdownRendererRightBrace{} \markdownRendererBackslash{}cdots p\markdownRendererUnderscore{}k\markdownRendererCircumflex{}\markdownRendererLeftBrace{}\markdownRendererBackslash{}alpha\markdownRendererUnderscore{}k\markdownRendererRightBrace{}\markdownRendererSoftLineBreak
{}\markdownRendererDollarSign{}\markdownRendererDollarSign{}\markdownRendererParagraphSeparator
{}\markdownRendererDollarSign{}\markdownRendererBackslash{}clearpage\markdownRendererDollarSign{}\markdownRendererInterblockSeparator
{}\markdownRendererSectionBegin
\markdownRendererHeadingThree{Cantidad de factores}\markdownRendererInterblockSeparator
{}Sea \markdownRendererDollarSign{}\markdownRendererBackslash{}tau(n)\markdownRendererDollarSign{} el número de factores de un entero n. Por ejemplo, \markdownRendererDollarSign{}\markdownRendererBackslash{}tau(12) = 6\markdownRendererDollarSign{}, porque los factores de 12 son 1, 2, 3, 4, 6 y 12. Para calcular el valor de \markdownRendererDollarSign{}\markdownRendererBackslash{}tau(n)\markdownRendererDollarSign{}, podemos usar la fórmula\markdownRendererParagraphSeparator
{}\markdownRendererDollarSign{}\markdownRendererDollarSign{}\markdownRendererSoftLineBreak
{}\markdownRendererBackslash{}tau(n) = \markdownRendererBackslash{}prod\markdownRendererUnderscore{}\markdownRendererLeftBrace{}i=1\markdownRendererRightBrace{}\markdownRendererCircumflex{}\markdownRendererLeftBrace{}k\markdownRendererRightBrace{} (\markdownRendererBackslash{}alpha\markdownRendererUnderscore{}i + 1)\markdownRendererSoftLineBreak
{}\markdownRendererDollarSign{}\markdownRendererDollarSign{}\markdownRendererParagraphSeparator
{}ya que para cada primo \markdownRendererDollarSign{}p\markdownRendererUnderscore{}i\markdownRendererDollarSign{}, hay \markdownRendererDollarSign{}\markdownRendererBackslash{}alpha\markdownRendererUnderscore{}i + 1)\markdownRendererDollarSign{} formas de elegir cuántas veces aparece en el factor. Por ejemplo, como \markdownRendererDollarSign{}12 = 2\markdownRendererCircumflex{}2 \markdownRendererBackslash{}cdot 3\markdownRendererCircumflex{}1\markdownRendererDollarSign{} \markdownRendererDollarSign{}\markdownRendererBackslash{}tau(12) = (2+1) \markdownRendererBackslash{}cdot 1(2) = 6\markdownRendererDollarSign{}.\markdownRendererParagraphSeparator
{}\markdownRendererStrongEmphasis{Complejidad:} \markdownRendererDollarSign{}O(\markdownRendererBackslash{}sqrt\markdownRendererLeftBrace{}n\markdownRendererRightBrace{})\markdownRendererDollarSign{}\markdownRendererInterblockSeparator
{}\markdownRendererInputFencedCode{./_markdown_main/0988881d17522e0e94a93016be679e15.verbatim}{cpp}{cpp}\markdownRendererInterblockSeparator
{}\markdownRendererDollarSign{}\markdownRendererBackslash{}clearpage\markdownRendererDollarSign{}\markdownRendererInterblockSeparator
{}
\markdownRendererSectionEnd 
\markdownRendererSectionEnd \markdownRendererSectionBegin
\markdownRendererHeadingTwo{Suma de factores}\markdownRendererInterblockSeparator
{}Sea \markdownRendererDollarSign{}\markdownRendererBackslash{}sigma(n)\markdownRendererDollarSign{} la suma de los factores de un entero n. Por ejemplo, \markdownRendererDollarSign{}\markdownRendererBackslash{}sigma(12) = 28\markdownRendererDollarSign{}, porque (1 + 2 + 3 + 4 + 6 + 12 = 28). Para calcular el valor de \markdownRendererDollarSign{}\markdownRendererBackslash{}sigma(n)\markdownRendererDollarSign{}, podemos usar la fórmula\markdownRendererParagraphSeparator
{}\markdownRendererDollarSign{}\markdownRendererDollarSign{}\markdownRendererSoftLineBreak
{}\markdownRendererBackslash{}sigma(n) = \markdownRendererBackslash{}prod\markdownRendererUnderscore{}\markdownRendererLeftBrace{}i=1\markdownRendererRightBrace{}\markdownRendererCircumflex{}\markdownRendererLeftBrace{}k\markdownRendererRightBrace{} (1 + p\markdownRendererUnderscore{}i + \markdownRendererBackslash{}dots + p\markdownRendererUnderscore{}i\markdownRendererCircumflex{}\markdownRendererLeftBrace{}\markdownRendererBackslash{}alpha\markdownRendererUnderscore{}i\markdownRendererRightBrace{}) = \markdownRendererBackslash{}prod\markdownRendererUnderscore{}\markdownRendererLeftBrace{}i=1\markdownRendererRightBrace{}\markdownRendererCircumflex{}\markdownRendererLeftBrace{}k\markdownRendererRightBrace{} \markdownRendererBackslash{}frac\markdownRendererLeftBrace{}p\markdownRendererUnderscore{}i\markdownRendererCircumflex{}\markdownRendererLeftBrace{}\markdownRendererBackslash{}alpha\markdownRendererUnderscore{}i+1\markdownRendererRightBrace{} - 1\markdownRendererRightBrace{}\markdownRendererLeftBrace{}p\markdownRendererUnderscore{}i - 1\markdownRendererRightBrace{}\markdownRendererSoftLineBreak
{}\markdownRendererDollarSign{}\markdownRendererDollarSign{}\markdownRendererParagraphSeparator
{}donde la última forma se basa en la fórmula de progresión geométrica.\markdownRendererSoftLineBreak
{}Por ejemplo, \markdownRendererDollarSign{} \markdownRendererBackslash{}sigma(12) = \markdownRendererBackslash{}frac\markdownRendererLeftBrace{}2\markdownRendererCircumflex{}3 - 1\markdownRendererRightBrace{}\markdownRendererLeftBrace{}2 - 1\markdownRendererRightBrace{} \markdownRendererBackslash{}cdot \markdownRendererBackslash{}frac\markdownRendererLeftBrace{}3\markdownRendererCircumflex{}2 - 1\markdownRendererRightBrace{}\markdownRendererLeftBrace{}3 - 1\markdownRendererRightBrace{} = 28\markdownRendererDollarSign{}\markdownRendererParagraphSeparator
{}\markdownRendererStrongEmphasis{Complejidad:} \markdownRendererDollarSign{}O(\markdownRendererBackslash{}sqrt\markdownRendererLeftBrace{}n\markdownRendererRightBrace{})\markdownRendererDollarSign{}\markdownRendererInterblockSeparator
{}\markdownRendererInputFencedCode{./_markdown_main/306b0cf2d68e968c0c2942f028c8ca92.verbatim}{cpp}{cpp}\markdownRendererInterblockSeparator
{}\markdownRendererDollarSign{}\markdownRendererBackslash{}clearpage\markdownRendererDollarSign{}\markdownRendererInterblockSeparator
{}\markdownRendererSectionBegin
\markdownRendererHeadingThree{Lista de factores}\markdownRendererInterblockSeparator
{}Observa que cada factor primo aparece en el vector tantas veces como divide el número.\markdownRendererParagraphSeparator
{}Por ejemplo \markdownRendererDollarSign{}12 = 2\markdownRendererCircumflex{}2 * 3\markdownRendererDollarSign{}, por lo que el resultado de la función es \markdownRendererDollarSign{}[2,2,3]\markdownRendererDollarSign{}\markdownRendererParagraphSeparator
{}\markdownRendererStrongEmphasis{Complejidad:} \markdownRendererDollarSign{}O(\markdownRendererBackslash{}sqrt\markdownRendererLeftBrace{}n\markdownRendererRightBrace{})\markdownRendererDollarSign{}\markdownRendererInterblockSeparator
{}\markdownRendererInputFencedCode{./_markdown_main/c2c22493926f2ec28a6855dd502f2644.verbatim}{cpp}{cpp}\markdownRendererInterblockSeparator
{}Variante que toma la criba y devuelve un map donde el valor de cada clave es la potencia, ejemplo\markdownRendererSoftLineBreak
{}\markdownRendererDollarSign{}12 = \markdownRendererLeftBrace{}2:2, 3:1\markdownRendererRightBrace{}\markdownRendererDollarSign{}\markdownRendererParagraphSeparator
{}\markdownRendererStrongEmphasis{Complejidad:} \markdownRendererDollarSign{}O(\markdownRendererBackslash{}pi(\markdownRendererBackslash{}sqrt\markdownRendererLeftBrace{}n\markdownRendererRightBrace{}) + \markdownRendererBackslash{}log n) \markdownRendererBackslash{}approx O(\markdownRendererBackslash{}frac\markdownRendererLeftBrace{}\markdownRendererBackslash{}sqrt\markdownRendererLeftBrace{}n\markdownRendererRightBrace{}\markdownRendererRightBrace{}\markdownRendererLeftBrace{}\markdownRendererBackslash{}log n\markdownRendererRightBrace{} + \markdownRendererBackslash{}log n)\markdownRendererDollarSign{}\markdownRendererInterblockSeparator
{}\markdownRendererInputFencedCode{./_markdown_main/cc0d3c823343b546c86096ee4df89adc.verbatim}{cpp}{cpp}\markdownRendererInterblockSeparator
{}
\markdownRendererSectionEnd 
\markdownRendererSectionEnd \markdownRendererSectionBegin
\markdownRendererHeadingTwo{Pruebas de primalidad}\markdownRendererInterblockSeparator
{}Revisar si un número es primo.\markdownRendererInterblockSeparator
{}\markdownRendererInputFencedCode{./_markdown_main/97dcb86dfa1386951a76c12510d6f0ad.verbatim}{cpp}{cpp}\markdownRendererInterblockSeparator
{}\markdownRendererSectionBegin
\markdownRendererHeadingThree{Algoritmos probabilisticos}\markdownRendererInterblockSeparator
{}Estos algorimtos se basan en comprobar de manera probabilistica basado en el Teoerema debil de Fermat, cuanto mayor sea la cantidad de iteraciones, mneor es la probabilidad de equviocarse.\markdownRendererInterblockSeparator
{}\markdownRendererSectionBegin
\markdownRendererHeadingFour{Prueba de Fermat}\markdownRendererInterblockSeparator
{}La prueba de Fermat puede fallar con los números de Charmichael  (\markdownRendererDollarSign{}a\markdownRendererCircumflex{}\markdownRendererLeftBrace{}n-1\markdownRendererRightBrace{} \markdownRendererBackslash{}equiv 1 \markdownRendererBackslash{}bmod n\markdownRendererDollarSign{} se cumple para todo\markdownRendererDollarSign{}a\markdownRendererDollarSign{} coprimo con \markdownRendererDollarSign{}n\markdownRendererDollarSign{}, por ejemplo \markdownRendererDollarSign{}561 = 3 \markdownRendererBackslash{}times 11 \markdownRendererBackslash{}times 17\markdownRendererDollarSign{}), sin embargo en la practica se sigue usando ya que es bastante rápida y los números de Charmichael son raros.\markdownRendererInterblockSeparator
{}\markdownRendererInputFencedCode{./_markdown_main/29fa9c22c56ad85eaa06176f9d7aeec6.verbatim}{cpp}{cpp}\markdownRendererInterblockSeparator
{}
\markdownRendererSectionEnd \markdownRendererSectionBegin
\markdownRendererHeadingFour{Prueba de Miller}\markdownRendererInterblockSeparator
{}Miller desmostro que el algoritmo se puede volver determinista si se prueban todas bases \markdownRendererDollarSign{}a \markdownRendererBackslash{}leq 2 ln(n)\markdownRendererCircumflex{}2\markdownRendererDollarSign{},\markdownRendererSoftLineBreak
{}siendo que para los números de 64 bit, basta con revisar las primeras 12 bases primas.\markdownRendererParagraphSeparator
{}\markdownRendererInputFencedCode{./_markdown_main/5942cc6a6c9d379f465331b341faa41d.verbatim}{cpp}{cpp}\markdownRendererInterblockSeparator
{}\markdownRendererDollarSign{}\markdownRendererBackslash{}clearpage\markdownRendererDollarSign{}\markdownRendererInterblockSeparator
{}
\markdownRendererSectionEnd 
\markdownRendererSectionEnd \markdownRendererSectionBegin
\markdownRendererHeadingThree{Primos como diferencia de cuadrados}\markdownRendererParagraphSeparator
{}Un numero impar se puede escribir \markdownRendererDollarSign{}n = pq\markdownRendererDollarSign{} como la diferencia de cuadrados \markdownRendererDollarSign{}n = a\markdownRendererCircumflex{}2 - b\markdownRendererCircumflex{}2\markdownRendererDollarSign{}.\markdownRendererParagraphSeparator
{}\markdownRendererDollarSign{}\markdownRendererDollarSign{}n = (\markdownRendererBackslash{}frac\markdownRendererLeftBrace{}p+q\markdownRendererRightBrace{}\markdownRendererLeftBrace{}2\markdownRendererRightBrace{})\markdownRendererCircumflex{}2 - (\markdownRendererBackslash{}frac\markdownRendererLeftBrace{}p-q\markdownRendererRightBrace{}\markdownRendererLeftBrace{}2\markdownRendererRightBrace{})\markdownRendererCircumflex{}2\markdownRendererDollarSign{}\markdownRendererDollarSign{}\markdownRendererParagraphSeparator
{}Estos valores \markdownRendererDollarSign{}p\markdownRendererDollarSign{} y \markdownRendererDollarSign{}q\markdownRendererDollarSign{} se pueden obtener con el metodo de factorización de Fermat. \markdownRendererDollarSign{}O(\markdownRendererPipe{}p-q\markdownRendererPipe{})\markdownRendererDollarSign{} cabe\markdownRendererSoftLineBreak
{}mencionar que si los números son demasiado grande o la distancia entre \markdownRendererDollarSign{}p\markdownRendererDollarSign{} y \markdownRendererDollarSign{}q\markdownRendererDollarSign{}, este algoritmo es\markdownRendererSoftLineBreak
{}extremadamente lento.\markdownRendererInterblockSeparator
{}\markdownRendererInputFencedCode{./_markdown_main/a9a4b39bbf10fd16685c491bd77951e3.verbatim}{cpp}{cpp}\markdownRendererInterblockSeparator
{}
\markdownRendererSectionEnd 
\markdownRendererSectionEnd \markdownRendererSectionBegin
\markdownRendererHeadingTwo{Computar números primos}\markdownRendererParagraphSeparator
{}\markdownRendererSectionBegin
\markdownRendererHeadingThree{Criba de Erastotenes}\markdownRendererInterblockSeparator
{}\markdownRendererStrongEmphasis{Idea:} marcar múltiplos de cada primo hasta \markdownRendererDollarSign{}n\markdownRendererDollarSign{}.\markdownRendererHardLineBreak
{}\markdownRendererStrongEmphasis{Complejidad:} \markdownRendererDollarSign{}O(n \markdownRendererBackslash{}log \markdownRendererBackslash{}log n)\markdownRendererDollarSign{}\markdownRendererInterblockSeparator
{}\markdownRendererInputFencedCode{./_markdown_main/62198306d6a4f3c98cb7c8f3275e0616.verbatim}{cpp}{cpp}\markdownRendererInterblockSeparator
{}\markdownRendererDollarSign{}\markdownRendererBackslash{}clearpage\markdownRendererDollarSign{}\markdownRendererInterblockSeparator
{}
\markdownRendererSectionEnd \markdownRendererSectionBegin
\markdownRendererHeadingThree{Criba Lineal}\markdownRendererInterblockSeparator
{}\markdownRendererStrongEmphasis{Idea:} Existe una optimización que permite realizar la criba en \markdownRendererDollarSign{}O(n)\markdownRendererDollarSign{} a cambio de usar más memoria, que solo debe ser usada si \markdownRendererDollarSign{}N \markdownRendererBackslash{}leq 10\markdownRendererCircumflex{}7\markdownRendererDollarSign{}\markdownRendererSoftLineBreak
{}\markdownRendererStrongEmphasis{Complejidad} : \markdownRendererDollarSign{}O(n)\markdownRendererDollarSign{}\markdownRendererInterblockSeparator
{}\markdownRendererInputFencedCode{./_markdown_main/b9d19f4fd13f5a17a42e8a28995b2ec1.verbatim}{cpp}{cpp}\markdownRendererInterblockSeparator
{}
\markdownRendererSectionEnd \markdownRendererSectionBegin
\markdownRendererHeadingThree{Contar número de primos en un rango}\markdownRendererInterblockSeparator
{}\markdownRendererStrongEmphasis{Idea:} Realizar una suma prefijos\markdownRendererSoftLineBreak
{}\markdownRendererStrongEmphasis{Complejidad de precomputo:} \markdownRendererDollarSign{}O(n \markdownRendererBackslash{}log \markdownRendererBackslash{}log \markdownRendererBackslash{}sqrt\markdownRendererLeftBrace{}n\markdownRendererRightBrace{})\markdownRendererDollarSign{}\markdownRendererSoftLineBreak
{}\markdownRendererStrongEmphasis{Complejidad de consulta:} \markdownRendererDollarSign{}O(1)\markdownRendererDollarSign{}\markdownRendererInterblockSeparator
{}\markdownRendererInputFencedCode{./_markdown_main/564d090d7d9420c98260c0d95094053f.verbatim}{cpp}{cpp}\markdownRendererInterblockSeparator
{}\markdownRendererDollarSign{}\markdownRendererBackslash{}clearpage\markdownRendererDollarSign{}\markdownRendererInterblockSeparator
{}
\markdownRendererSectionEnd \markdownRendererSectionBegin
\markdownRendererHeadingThree{Encontrar primos en un rango}\markdownRendererInterblockSeparator
{}Encontrar primos en un rango\markdownRendererParagraphSeparator
{}Encontrar números primos en un rango \markdownRendererDollarSign{}[L, R]\markdownRendererDollarSign{} donde \markdownRendererDollarSign{}R - L + 1 \markdownRendererBackslash{}approx 10\markdownRendererCircumflex{}7\markdownRendererDollarSign{} y \markdownRendererDollarSign{}R \markdownRendererBackslash{}leq 10\markdownRendererCircumflex{}\markdownRendererLeftBrace{}12\markdownRendererRightBrace{}\markdownRendererDollarSign{}.\markdownRendererInterblockSeparator
{}\markdownRendererInputFencedCode{./_markdown_main/68cbc9644aa9ef02e1213c2627a4fd4e.verbatim}{cpp}{cpp}\markdownRendererInterblockSeparator
{}
\markdownRendererSectionEnd 
\markdownRendererSectionEnd \markdownRendererSectionBegin
\markdownRendererHeadingTwo{Algoritmo de Euclides}\markdownRendererParagraphSeparator
{}\markdownRendererSectionBegin
\markdownRendererHeadingThree{Máximo Comúm Divisior}\markdownRendererInterblockSeparator
{}\markdownRendererDollarSign{}\markdownRendererDollarSign{}\markdownRendererSoftLineBreak
{}\markdownRendererBackslash{}gcd(a,b) = \markdownRendererBackslash{}gcd(b, a \markdownRendererBackslash{}bmod b)\markdownRendererSoftLineBreak
{}\markdownRendererDollarSign{}\markdownRendererDollarSign{}\markdownRendererParagraphSeparator
{}\markdownRendererStrongEmphasis{Complejidad:} \markdownRendererDollarSign{}O(\markdownRendererBackslash{}log \markdownRendererBackslash{}min(a,b))\markdownRendererDollarSign{}.\markdownRendererInterblockSeparator
{}\markdownRendererInputFencedCode{./_markdown_main/e5a7f9b9c711a56f4dd40c74d257f0e2.verbatim}{c}{c}\markdownRendererInterblockSeparator
{}
\markdownRendererSectionEnd \markdownRendererSectionBegin
\markdownRendererHeadingThree{Máximo Común Multiplo}\markdownRendererInterblockSeparator
{}\markdownRendererDollarSign{}\markdownRendererDollarSign{}\markdownRendererSoftLineBreak
{}\markdownRendererBackslash{}mathrm\markdownRendererLeftBrace{}lcm\markdownRendererRightBrace{}(a,b) = \markdownRendererBackslash{}frac\markdownRendererLeftBrace{}a \markdownRendererBackslash{}cdot b\markdownRendererRightBrace{}\markdownRendererLeftBrace{}\markdownRendererBackslash{}gcd(a,b)\markdownRendererRightBrace{}\markdownRendererSoftLineBreak
{}\markdownRendererDollarSign{}\markdownRendererDollarSign{}\markdownRendererParagraphSeparator
{}\markdownRendererStrongEmphasis{Complejidad:} \markdownRendererDollarSign{}O(\markdownRendererBackslash{}log \markdownRendererBackslash{}min(a,b))\markdownRendererDollarSign{}.\markdownRendererInterblockSeparator
{}\markdownRendererInputFencedCode{./_markdown_main/394daf27a502d67ee061bed27f75b198.verbatim}{c}{c}\markdownRendererInterblockSeparator
{}\markdownRendererDollarSign{}\markdownRendererBackslash{}clearpage\markdownRendererDollarSign{}\markdownRendererInterblockSeparator
{}
\markdownRendererSectionEnd \markdownRendererSectionBegin
\markdownRendererHeadingThree{Euclides Extendido}\markdownRendererInterblockSeparator
{}\markdownRendererStrongEmphasis{Devuelve} \markdownRendererDollarSign{}x,y\markdownRendererDollarSign{} tales que \markdownRendererDollarSign{}a x + b y = \markdownRendererBackslash{}gcd(a,b)\markdownRendererDollarSign{} — útil para inversos modulares y diofánticas.\markdownRendererInterblockSeparator
{}\markdownRendererInputFencedCode{./_markdown_main/6534d8ec2fb7c3e13807298b70eadf5a.verbatim}{cpp}{cpp}\markdownRendererInterblockSeparator
{}
\markdownRendererSectionEnd \markdownRendererSectionBegin
\markdownRendererHeadingThree{Numeros Coprimos}\markdownRendererInterblockSeparator
{}Dos enteros \markdownRendererDollarSign{}a\markdownRendererDollarSign{} y \markdownRendererDollarSign{}b\markdownRendererDollarSign{} son coprimos si \markdownRendererDollarSign{}\markdownRendererBackslash{}gcd(a,b) = 1\markdownRendererDollarSign{}.\markdownRendererInterblockSeparator
{}\markdownRendererInputFencedCode{./_markdown_main/62ab9e372f52b49dbcba35952dccf31d.verbatim}{cpp}{cpp}\markdownRendererInterblockSeparator
{}
\markdownRendererSectionEnd \markdownRendererSectionBegin
\markdownRendererHeadingThree{Funcion Totiente de Euler}\markdownRendererInterblockSeparator
{}La \markdownRendererStrongEmphasis{función totiente de Euler} \markdownRendererDollarSign{}\markdownRendererBackslash{}varphi(n)\markdownRendererDollarSign{} da el número de enteros entre \markdownRendererDollarSign{}1\markdownRendererDollarSign{} y \markdownRendererDollarSign{}n\markdownRendererDollarSign{} que son coprimos con \markdownRendererDollarSign{}n\markdownRendererDollarSign{} si\markdownRendererSoftLineBreak
{}\markdownRendererDollarSign{}\markdownRendererDollarSign{}\markdownRendererSoftLineBreak
{}n = p\markdownRendererUnderscore{}1\markdownRendererCircumflex{}\markdownRendererLeftBrace{}\markdownRendererBackslash{}alpha\markdownRendererUnderscore{}1\markdownRendererRightBrace{} \markdownRendererBackslash{}cdots p\markdownRendererUnderscore{}k\markdownRendererCircumflex{}\markdownRendererLeftBrace{}\markdownRendererBackslash{}alpha\markdownRendererUnderscore{}k\markdownRendererRightBrace{},\markdownRendererSoftLineBreak
{}\markdownRendererDollarSign{}\markdownRendererDollarSign{}\markdownRendererParagraphSeparator
{}\markdownRendererDollarSign{}\markdownRendererDollarSign{}\markdownRendererSoftLineBreak
{}\markdownRendererBackslash{}varphi(n) = \markdownRendererBackslash{}prod\markdownRendererUnderscore{}\markdownRendererLeftBrace{}i=1\markdownRendererRightBrace{}\markdownRendererCircumflex{}\markdownRendererLeftBrace{}k\markdownRendererRightBrace{} p\markdownRendererUnderscore{}i\markdownRendererCircumflex{}\markdownRendererLeftBrace{}\markdownRendererBackslash{}alpha\markdownRendererUnderscore{}i - 1\markdownRendererRightBrace{} (p\markdownRendererUnderscore{}i - 1)\markdownRendererSoftLineBreak
{}\markdownRendererDollarSign{}\markdownRendererDollarSign{}\markdownRendererParagraphSeparator
{}\markdownRendererStrongEmphasis{Ejemplo:} \markdownRendererDollarSign{}10 = 2 \markdownRendererBackslash{}cdot 5 ;\markdownRendererBackslash{}Rightarrow; \markdownRendererBackslash{}varphi(10) = 2\markdownRendererCircumflex{}\markdownRendererLeftBrace{}1-1\markdownRendererRightBrace{}(2-1) \markdownRendererBackslash{}cdot 5\markdownRendererCircumflex{}\markdownRendererLeftBrace{}1-1\markdownRendererRightBrace{}(5-1) = 4\markdownRendererDollarSign{}\markdownRendererParagraphSeparator
{}\markdownRendererStrongEmphasis{Complejidad:} \markdownRendererDollarSign{}O(\markdownRendererBackslash{}sqrt\markdownRendererLeftBrace{}n\markdownRendererRightBrace{})\markdownRendererDollarSign{}\markdownRendererInterblockSeparator
{}\markdownRendererInputFencedCode{./_markdown_main/ce2a36645f71b493c7aa5aa988ec103c.verbatim}{cpp}{cpp}\markdownRendererInterblockSeparator
{}\markdownRendererDollarSign{}\markdownRendererBackslash{}clearpage\markdownRendererDollarSign{}\markdownRendererInterblockSeparator
{}
\markdownRendererSectionEnd \markdownRendererSectionBegin
\markdownRendererHeadingThree{Teorema de Euler}\markdownRendererInterblockSeparator
{}\markdownRendererStrongEmphasis{Teorema de Euler:}\markdownRendererParagraphSeparator
{}\markdownRendererDollarSign{}\markdownRendererDollarSign{}\markdownRendererSoftLineBreak
{}x\markdownRendererCircumflex{}\markdownRendererLeftBrace{}\markdownRendererBackslash{}varphi(m)\markdownRendererRightBrace{} \markdownRendererBackslash{}equiv 1 \markdownRendererBackslash{}pmod\markdownRendererLeftBrace{}m\markdownRendererRightBrace{}, \markdownRendererBackslash{}quad x \markdownRendererBackslash{}text\markdownRendererLeftBrace{} coprimo con \markdownRendererRightBrace{} m\markdownRendererSoftLineBreak
{}\markdownRendererDollarSign{}\markdownRendererDollarSign{}\markdownRendererParagraphSeparator
{}Si \markdownRendererDollarSign{}m\markdownRendererDollarSign{} es primo (\markdownRendererDollarSign{}\markdownRendererBackslash{}varphi(m) = m - 1\markdownRendererDollarSign{}):\markdownRendererInterblockSeparator
{}
\markdownRendererSectionEnd \markdownRendererSectionBegin
\markdownRendererHeadingThree{Pequeño Teorema de Fermat}\markdownRendererInterblockSeparator
{}\markdownRendererDollarSign{}\markdownRendererDollarSign{}\markdownRendererSoftLineBreak
{}x\markdownRendererCircumflex{}\markdownRendererLeftBrace{}m - 1\markdownRendererRightBrace{} \markdownRendererBackslash{}equiv 1 \markdownRendererBackslash{}pmod\markdownRendererLeftBrace{}m\markdownRendererRightBrace{}\markdownRendererSoftLineBreak
{}\markdownRendererDollarSign{}\markdownRendererDollarSign{}\markdownRendererParagraphSeparator
{}
\markdownRendererSectionEnd 
\markdownRendererSectionEnd \markdownRendererSectionBegin
\markdownRendererHeadingTwo{Inversos Multiplicativos Modulares}\markdownRendererInterblockSeparator
{}El inverso multiplicativo de \markdownRendererDollarSign{}x\markdownRendererDollarSign{} módulo \markdownRendererDollarSign{}m\markdownRendererDollarSign{}, \markdownRendererDollarSign{}\markdownRendererBackslash{}mathrm\markdownRendererLeftBrace{}inv\markdownRendererRightBrace{}\markdownRendererUnderscore{}m(x)\markdownRendererDollarSign{}, cumple:\markdownRendererParagraphSeparator
{}\markdownRendererDollarSign{}\markdownRendererDollarSign{}\markdownRendererSoftLineBreak
{}x \markdownRendererBackslash{}cdot \markdownRendererBackslash{}mathrm\markdownRendererLeftBrace{}inv\markdownRendererRightBrace{}\markdownRendererUnderscore{}m(x) \markdownRendererBackslash{}equiv 1 \markdownRendererBackslash{}pmod\markdownRendererLeftBrace{}m\markdownRendererRightBrace{}\markdownRendererSoftLineBreak
{}\markdownRendererDollarSign{}\markdownRendererDollarSign{}\markdownRendererParagraphSeparator
{}Existe si y solo si \markdownRendererDollarSign{}\markdownRendererBackslash{}gcd(x, m) = 1\markdownRendererDollarSign{}, y se puede calcular con:\markdownRendererParagraphSeparator
{}\markdownRendererDollarSign{}\markdownRendererDollarSign{}\markdownRendererSoftLineBreak
{}\markdownRendererBackslash{}mathrm\markdownRendererLeftBrace{}inv\markdownRendererRightBrace{}\markdownRendererUnderscore{}m(x) = x\markdownRendererCircumflex{}\markdownRendererLeftBrace{}\markdownRendererBackslash{}varphi(m) - 1\markdownRendererRightBrace{} \markdownRendererBackslash{}pmod\markdownRendererLeftBrace{}m\markdownRendererRightBrace{}\markdownRendererSoftLineBreak
{}\markdownRendererDollarSign{}\markdownRendererDollarSign{}\markdownRendererParagraphSeparator
{}Si \markdownRendererDollarSign{}m\markdownRendererDollarSign{} es primo:\markdownRendererParagraphSeparator
{}\markdownRendererDollarSign{}\markdownRendererDollarSign{}\markdownRendererSoftLineBreak
{}\markdownRendererBackslash{}mathrm\markdownRendererLeftBrace{}inv\markdownRendererRightBrace{}\markdownRendererUnderscore{}m(x) = x\markdownRendererCircumflex{}\markdownRendererLeftBrace{}m - 2\markdownRendererRightBrace{} \markdownRendererBackslash{}pmod\markdownRendererLeftBrace{}m\markdownRendererRightBrace{}\markdownRendererSoftLineBreak
{}\markdownRendererDollarSign{}\markdownRendererDollarSign{}\markdownRendererParagraphSeparator
{}\markdownRendererStrongEmphasis{Ejemplo:} \markdownRendererDollarSign{}\markdownRendererBackslash{}mathrm\markdownRendererLeftBrace{}inv\markdownRendererRightBrace{}\markdownRendererUnderscore{}\markdownRendererLeftBrace{}17\markdownRendererRightBrace{}(6) = 6\markdownRendererCircumflex{}\markdownRendererLeftBrace{}15\markdownRendererRightBrace{} \markdownRendererBackslash{}bmod 17 = 3\markdownRendererDollarSign{}\markdownRendererParagraphSeparator
{}\markdownRendererInputFencedCode{./_markdown_main/ef741a8ede8e19e8bfca2376eb0ca5ae.verbatim}{cpp}{cpp}\markdownRendererInterblockSeparator
{}\markdownRendererDollarSign{}\markdownRendererBackslash{}clearpage\markdownRendererDollarSign{}\markdownRendererInterblockSeparator
{}
\markdownRendererSectionEnd \markdownRendererSectionBegin
\markdownRendererHeadingTwo{Ecuaciones Diofánticas}\markdownRendererInterblockSeparator
{}Una ecuación diofántica tiene la forma:\markdownRendererParagraphSeparator
{}\markdownRendererDollarSign{}\markdownRendererDollarSign{}\markdownRendererSoftLineBreak
{}a x + b y = c, \markdownRendererBackslash{}quad a, b, c \markdownRendererBackslash{}in \markdownRendererBackslash{}mathbb\markdownRendererLeftBrace{}Z\markdownRendererRightBrace{}\markdownRendererSoftLineBreak
{}\markdownRendererDollarSign{}\markdownRendererDollarSign{}\markdownRendererParagraphSeparator
{}Se puede resolver si \markdownRendererDollarSign{}c\markdownRendererDollarSign{} es divisible por \markdownRendererDollarSign{}\markdownRendererBackslash{}gcd(a,b)\markdownRendererDollarSign{}.\markdownRendererHardLineBreak
{}Si \markdownRendererDollarSign{}(x\markdownRendererUnderscore{}0, y\markdownRendererUnderscore{}0)\markdownRendererDollarSign{} es una solución, todas las soluciones son:\markdownRendererParagraphSeparator
{}\markdownRendererDollarSign{}\markdownRendererDollarSign{}\markdownRendererSoftLineBreak
{}\markdownRendererBackslash{}begin\markdownRendererLeftBrace{}cases\markdownRendererRightBrace{}\markdownRendererSoftLineBreak
{}x = x\markdownRendererUnderscore{}0 + k \markdownRendererBackslash{}dfrac\markdownRendererLeftBrace{}b\markdownRendererRightBrace{}\markdownRendererLeftBrace{}\markdownRendererBackslash{}gcd(a,b)\markdownRendererRightBrace{} \markdownRendererBackslash{}\markdownRendererBackslash{}\markdownRendererSoftLineBreak
{}y = y\markdownRendererUnderscore{}0 - k \markdownRendererBackslash{}dfrac\markdownRendererLeftBrace{}a\markdownRendererRightBrace{}\markdownRendererLeftBrace{}\markdownRendererBackslash{}gcd(a,b)\markdownRendererRightBrace{} \markdownRendererBackslash{}\markdownRendererBackslash{}\markdownRendererSoftLineBreak
{}k \markdownRendererBackslash{}in \markdownRendererBackslash{}mathbb\markdownRendererLeftBrace{}Z\markdownRendererRightBrace{}\markdownRendererSoftLineBreak
{}\markdownRendererBackslash{}end\markdownRendererLeftBrace{}cases\markdownRendererRightBrace{}\markdownRendererSoftLineBreak
{}\markdownRendererDollarSign{}\markdownRendererDollarSign{}\markdownRendererParagraphSeparator
{}\markdownRendererStrongEmphasis{Ejemplo:} resolver \markdownRendererDollarSign{}39x + 15y = 12\markdownRendererDollarSign{} \markdownRendererDollarSign{}\markdownRendererBackslash{}gcd(39, 15) = 3 \markdownRendererBackslash{}mid 12\markdownRendererDollarSign{}\markdownRendererParagraphSeparator
{}Usando el algoritmo de Euclides extendido:\markdownRendererParagraphSeparator
{}\markdownRendererDollarSign{}39 \markdownRendererBackslash{}cdot 2 + 15 \markdownRendererBackslash{}cdot (-5) = 3 \markdownRendererBackslash{}implies 39 \markdownRendererBackslash{}cdot 8 + 15 \markdownRendererBackslash{}cdot (-20) = 12\markdownRendererDollarSign{}\markdownRendererInterblockSeparator
{}\markdownRendererInputFencedCode{./_markdown_main/22c17f3e7186de3b84533c4fc0500a3c.verbatim}{cpp}{cpp}\markdownRendererInterblockSeparator
{}\markdownRendererDollarSign{}\markdownRendererBackslash{}clearpage\markdownRendererDollarSign{}\markdownRendererInterblockSeparator
{}\markdownRendererSectionBegin
\markdownRendererHeadingThree{Teorema Chino del Resto}\markdownRendererInterblockSeparator
{}\markdownRendererStrongEmphasis{Sistema de congruencias:}\markdownRendererParagraphSeparator
{}\markdownRendererDollarSign{}\markdownRendererDollarSign{}\markdownRendererSoftLineBreak
{}\markdownRendererBackslash{}begin\markdownRendererLeftBrace{}cases\markdownRendererRightBrace{}\markdownRendererSoftLineBreak
{}x \markdownRendererBackslash{}equiv a\markdownRendererUnderscore{}1 \markdownRendererBackslash{}pmod\markdownRendererLeftBrace{}m\markdownRendererUnderscore{}1\markdownRendererRightBrace{} \markdownRendererBackslash{}\markdownRendererBackslash{}\markdownRendererSoftLineBreak
{}x \markdownRendererBackslash{}equiv a\markdownRendererUnderscore{}2 \markdownRendererBackslash{}pmod\markdownRendererLeftBrace{}m\markdownRendererUnderscore{}2\markdownRendererRightBrace{} \markdownRendererBackslash{}\markdownRendererBackslash{}\markdownRendererSoftLineBreak
{}\markdownRendererBackslash{}vdots \markdownRendererBackslash{}\markdownRendererBackslash{}\markdownRendererSoftLineBreak
{}x \markdownRendererBackslash{}equiv a\markdownRendererUnderscore{}n \markdownRendererBackslash{}pmod\markdownRendererLeftBrace{}m\markdownRendererUnderscore{}n\markdownRendererRightBrace{}\markdownRendererSoftLineBreak
{}\markdownRendererBackslash{}end\markdownRendererLeftBrace{}cases\markdownRendererRightBrace{}\markdownRendererSoftLineBreak
{}\markdownRendererDollarSign{}\markdownRendererDollarSign{}\markdownRendererParagraphSeparator
{}con \markdownRendererDollarSign{}m\markdownRendererUnderscore{}i\markdownRendererDollarSign{} coprimos dos a dos.\markdownRendererParagraphSeparator
{}\markdownRendererStrongEmphasis{Solución general:}\markdownRendererParagraphSeparator
{}\markdownRendererDollarSign{}\markdownRendererDollarSign{}\markdownRendererSoftLineBreak
{}x = \markdownRendererBackslash{}sum\markdownRendererUnderscore{}\markdownRendererLeftBrace{}k=1\markdownRendererRightBrace{}\markdownRendererCircumflex{}\markdownRendererLeftBrace{}n\markdownRendererRightBrace{} a\markdownRendererUnderscore{}k X\markdownRendererUnderscore{}k , \markdownRendererBackslash{}mathrm\markdownRendererLeftBrace{}inv\markdownRendererRightBrace{}\markdownRendererUnderscore{}\markdownRendererLeftBrace{}m\markdownRendererUnderscore{}k\markdownRendererRightBrace{}(X\markdownRendererUnderscore{}k)\markdownRendererSoftLineBreak
{}\markdownRendererDollarSign{}\markdownRendererDollarSign{}\markdownRendererParagraphSeparator
{}donde\markdownRendererParagraphSeparator
{}\markdownRendererDollarSign{}\markdownRendererDollarSign{}\markdownRendererSoftLineBreak
{}X\markdownRendererUnderscore{}k = \markdownRendererBackslash{}frac\markdownRendererLeftBrace{}m\markdownRendererUnderscore{}1 m\markdownRendererUnderscore{}2 \markdownRendererBackslash{}cdots m\markdownRendererUnderscore{}n\markdownRendererRightBrace{}\markdownRendererLeftBrace{}m\markdownRendererUnderscore{}k\markdownRendererRightBrace{}\markdownRendererSoftLineBreak
{}\markdownRendererDollarSign{}\markdownRendererDollarSign{}\markdownRendererParagraphSeparator
{}\markdownRendererStrongEmphasis{Ejemplo:}\markdownRendererParagraphSeparator
{}\markdownRendererDollarSign{}\markdownRendererDollarSign{}\markdownRendererSoftLineBreak
{}\markdownRendererBackslash{}begin\markdownRendererLeftBrace{}cases\markdownRendererRightBrace{}\markdownRendererSoftLineBreak
{}x \markdownRendererBackslash{}equiv 3 \markdownRendererBackslash{}pmod\markdownRendererLeftBrace{}5\markdownRendererRightBrace{} \markdownRendererBackslash{}\markdownRendererBackslash{}\markdownRendererSoftLineBreak
{}x \markdownRendererBackslash{}equiv 4 \markdownRendererBackslash{}pmod\markdownRendererLeftBrace{}7\markdownRendererRightBrace{} \markdownRendererBackslash{}\markdownRendererBackslash{}\markdownRendererSoftLineBreak
{}x \markdownRendererBackslash{}equiv 2 \markdownRendererBackslash{}pmod\markdownRendererLeftBrace{}3\markdownRendererRightBrace{}\markdownRendererSoftLineBreak
{}\markdownRendererBackslash{}end\markdownRendererLeftBrace{}cases\markdownRendererRightBrace{}\markdownRendererSoftLineBreak
{}\markdownRendererBackslash{}implies\markdownRendererSoftLineBreak
{}x = 263\markdownRendererSoftLineBreak
{}\markdownRendererDollarSign{}\markdownRendererDollarSign{}\markdownRendererParagraphSeparator
{}\markdownRendererStrongEmphasis{Infinitas soluciones:}\markdownRendererParagraphSeparator
{}\markdownRendererDollarSign{}\markdownRendererDollarSign{}\markdownRendererSoftLineBreak
{}x + m\markdownRendererUnderscore{}1 m\markdownRendererUnderscore{}2 \markdownRendererBackslash{}cdots m\markdownRendererUnderscore{}n\markdownRendererSoftLineBreak
{}\markdownRendererDollarSign{}\markdownRendererDollarSign{}\markdownRendererInterblockSeparator
{}\markdownRendererInputFencedCode{./_markdown_main/0c41cca03de1bd29536b2ce66e8c90c7.verbatim}{cpp}{cpp}\markdownRendererInterblockSeparator
{}
\markdownRendererSectionEnd 
\markdownRendererSectionEnd \markdownRendererSectionBegin
\markdownRendererHeadingTwo{MEX}\markdownRendererInterblockSeparator
{}El \markdownRendererStrongEmphasis{MEX} ("minimum excludant") de un conjunto de números enteros es el \markdownRendererStrongEmphasis{mínimo número entero no presente en el conjunto}.\markdownRendererParagraphSeparator
{}Ejemplo: \markdownRendererDollarSign{}\markdownRendererBackslash{}mathrm\markdownRendererLeftBrace{}mex\markdownRendererRightBrace{}(\markdownRendererLeftBrace{}0,1,3\markdownRendererRightBrace{}) = 2\markdownRendererDollarSign{}\markdownRendererInterblockSeparator
{}\markdownRendererInputFencedCode{./_markdown_main/35760a9625e11f9267eaf63a95669577.verbatim}{cpp}{cpp}\markdownRendererInterblockSeparator
{}
\markdownRendererSectionEnd 
\markdownRendererSectionEnd \markdownRendererSectionBegin
\markdownRendererHeadingOne{Juegos NIM}\markdownRendererInterblockSeparator
{}\markdownRendererUlBeginTight
\markdownRendererUlItem Hay (n) heaps con (x\markdownRendererUnderscore{}i) palillos.\markdownRendererUlItemEnd 
\markdownRendererUlItem Jugadores alternan turnos.\markdownRendererUlItemEnd 
\markdownRendererUlItem En cada turno se remueven palillos de un heap no vacío.\markdownRendererUlItemEnd 
\markdownRendererUlItem Gana quien remueve el último palillo.\markdownRendererUlItemEnd 
\markdownRendererUlEndTight \markdownRendererInterblockSeparator
{}\markdownRendererSectionBegin
\markdownRendererHeadingTwo{Nim Sum}\markdownRendererInterblockSeparator
{}\markdownRendererDollarSign{}\markdownRendererDollarSign{}\markdownRendererSoftLineBreak
{}s = x\markdownRendererUnderscore{}1 \markdownRendererBackslash{}oplus x\markdownRendererUnderscore{}2 \markdownRendererBackslash{}oplus \markdownRendererBackslash{}dots \markdownRendererBackslash{}oplus x\markdownRendererUnderscore{}n\markdownRendererSoftLineBreak
{}\markdownRendererDollarSign{}\markdownRendererDollarSign{}\markdownRendererInterblockSeparator
{}\markdownRendererUlBeginTight
\markdownRendererUlItem Si \markdownRendererDollarSign{}(s = 0)\markdownRendererDollarSign{} estado perdedor.\markdownRendererUlItemEnd 
\markdownRendererUlItem Si \markdownRendererDollarSign{}(s \markdownRendererBackslash{}neq 0)\markdownRendererDollarSign{} estado ganador.\markdownRendererUlItemEnd 
\markdownRendererUlEndTight \markdownRendererInterblockSeparator
{}
\markdownRendererSectionEnd \markdownRendererSectionBegin
\markdownRendererHeadingTwo{Movimiento óptimo}\markdownRendererInterblockSeparator
{}Para un estado ganador:\markdownRendererInterblockSeparator
{}\markdownRendererUlBeginTight
\markdownRendererUlItem Elegir heap \markdownRendererDollarSign{}(i)\markdownRendererDollarSign{} tal que \markdownRendererDollarSign{}(x\markdownRendererUnderscore{}i \markdownRendererBackslash{}oplus s < x\markdownRendererUnderscore{}i)\markdownRendererDollarSign{}\markdownRendererUlItemEnd 
\markdownRendererUlItem Remover hasta que el heap tenga \markdownRendererDollarSign{}(x\markdownRendererUnderscore{}i \markdownRendererBackslash{}oplus s)\markdownRendererDollarSign{} palillos.\markdownRendererUlItemEnd 
\markdownRendererUlEndTight \markdownRendererParagraphSeparator
{}
\markdownRendererSectionEnd \markdownRendererSectionBegin
\markdownRendererHeadingTwo{Teorema de Sprague–Grundy}\markdownRendererInterblockSeparator
{}El Teorema de Sprague–Grundy generaliza la estrategia de Nim a cualquier juego que cumpla:\markdownRendererInterblockSeparator
{}\markdownRendererUlBeginTight
\markdownRendererUlItem Dos jugadores alternan turnos.\markdownRendererUlItemEnd 
\markdownRendererUlItem El juego consiste en estados, y los movimientos posibles dependen solo del estado.\markdownRendererUlItemEnd 
\markdownRendererUlItem El juego termina cuando un jugador no puede moverse.\markdownRendererUlItemEnd 
\markdownRendererUlItem El juego termina seguro eventualmente.\markdownRendererUlItemEnd 
\markdownRendererUlItem Los jugadores tienen información completa y no hay aleatoriedad.\markdownRendererUlItemEnd 
\markdownRendererUlEndTight \markdownRendererInterblockSeparator
{}\markdownRendererDollarSign{}\markdownRendererBackslash{}clearpage\markdownRendererDollarSign{}\markdownRendererInterblockSeparator
{}
\markdownRendererSectionEnd \markdownRendererSectionBegin
\markdownRendererHeadingTwo{Números de Grundy}\markdownRendererInterblockSeparator
{}Cada estado del juego se asocia con un \markdownRendererStrongEmphasis{número de Grundy} (G(s)), equivalente al tamaño de un heap de Nim.\markdownRendererParagraphSeparator
{}Si \markdownRendererDollarSign{}(s)\markdownRendererDollarSign{} tiene movimientos hacia los estados \markdownRendererDollarSign{}(s\markdownRendererUnderscore{}1, s\markdownRendererUnderscore{}2, \markdownRendererBackslash{}dots, s\markdownRendererUnderscore{}n)\markdownRendererDollarSign{}:\markdownRendererParagraphSeparator
{}\markdownRendererDollarSign{}\markdownRendererDollarSign{}\markdownRendererSoftLineBreak
{}\markdownRendererBackslash{}text\markdownRendererLeftBrace{}mex\markdownRendererRightBrace{}(\markdownRendererLeftBrace{}g\markdownRendererUnderscore{}1,g\markdownRendererUnderscore{}2,\markdownRendererEllipsis{},g\markdownRendererUnderscore{}n\markdownRendererRightBrace{})\markdownRendererSoftLineBreak
{}\markdownRendererDollarSign{}\markdownRendererDollarSign{}
\markdownRendererSectionEnd 
\markdownRendererSectionEnd \markdownRendererDocumentEnd